\subsection{旋转方法}
\label{sec:ch4-method-of-rotations}

推论~\ref{cor:ch4-bounded-multiplier-kernel} 与 Plancherel 定理给出了 \eqref{eq:ch4-singular-integral-definition} 中的算子在 $L^2$ 上有界的充分条件. 本节与下一节发展 Calderón 和 Zygmund 的方法, 用以证明这些算子在 $L^p$, $1<p<\infty$, 上有界.

设 $T$ 是 $L^p(\mathbb R)$ 上的有界一维算子, 并取 $u\in S^{n-1}$. 令 $L_u=\{\lambda u:\lambda\in\mathbb R\}$, 以 $L_u^\perp$ 表示其正交补. 每个 $x\in\mathbb R^n$ 都唯一地写成 $x=x_1u+\bar x$, 其中 $x_1\in\mathbb R$, $\bar x\in L_u^\perp$. 定义 $T_uf(x)$ 为一维函数 $f(\,\cdot\,u+\bar x)$ 经 $T$ 作用后在 $x_1$ 处的值. 若 $C_p$ 是 $T$ 在 $L^p(\mathbb R)$ 上的范数, 则 Fubini 定理给出 $\|T_uf\|_p\le C_p\|f\|_p$.

\hypertarget{chapter-04-page-086}{}
以这种方式得到的算子包括方向 Hardy--Littlewood 极大函数
\[
 M_uf(x)=\sup_{h>0}\frac1{2h}\int_{-h}^h|f(x-tu)|\,dt,
\]
以及方向 Hilbert 变换
\[
 H_uf(x)=\frac1\pi\lim_{\varepsilon\to0}
 \int_{|t|>\varepsilon}f(x-tu)\frac{dt}{t}.
\]

由 $T$ 得到的算子 $T_u$ 在 $L^p(\mathbb R^n)$ 上一致有界, 因而它们的任何凸组合也有界. Minkowski 积分不等式立即给出下述结论.

\begin{Proposition}
\label{prop:ch4-directional-operator-average}
设一维算子 $T$ 在 $L^p(\mathbb R)$ 上有界, 范数为 $C_p$, 并令 $T_u$ 为由 $T$ 定义的方向算子. 对任意 $\Omega\in L^1(S^{n-1})$, 算子
\[
 T_\Omega f(x)=\int_{S^{n-1}}\Omega(u)T_uf(x)\,d\sigma(u)
\]
在 $L^p(\mathbb R^n)$ 上有界, 范数至多为 $C_p\|\Omega\|_1$.
\end{Proposition}

这个结果使我们能把一维结论提升到高维. 旋转方法最常见的应用是先采用极坐标, 再对径向部分使用方向算子. 例如, 对 $\Omega\in L^1(S^{n-1})$, 定义粗糙极大函数
\begin{equation}
\label{eq:ch4-rough-maximal-function}
 M_\Omega f(x)=\sup_{R>0}\frac1{|B(0,R)|}
 \int_{B(0,R)}|\Omega(y')|\,|f(x-y)|\,dy.
\end{equation}
把积分写成极坐标便得到
\[
 M_\Omega f(x)\le\frac1{|B(0,1)|}
 \int_{S^{n-1}}|\Omega(u)|M_uf(x)\,d\sigma(u).
\]

\hypertarget{chapter-04-page-087}{}
\begin{Corollary}
\label{cor:ch4-rough-maximal-lp}
若 $\Omega\in L^1(S^{n-1})$, 则 $M_\Omega$ 在 $L^p(\mathbb R^n)$, $1<p<\infty$, 上有界.
\end{Corollary}

若 $\Omega(u)=1$, 则 $M_\Omega$ 就是 $\mathbb R^n$ 中的 Hardy--Littlewood 极大函数. 旋转方法从一维情形推出它在 $L^p$, $p>1$, 上有界. 此方法给出的 $L^p$ 常数是 $O(n)$, 因为 $|S^{n-1}|/|B(0,1)|=n$(见第~\ref{subsec:ch2-size-of-constants}~节). 另一方面, 旋转方法不能给出弱型 $(1,1)$ 结论, 见第~\ref{subsec:ch4-weak-one-one}~节.

当 $\Omega$ 为奇函数时, 也可以把命题~\ref{prop:ch4-directional-operator-average} 用于 \eqref{eq:ch4-singular-integral-definition}. 固定 Schwartz 函数 $f$, 利用极坐标, $\Omega$ 的平均值为零以及控制收敛定理, 得
\[
 Tf(x)=\frac\pi2\int_{S^{n-1}}\Omega(u)H_uf(x)\,d\sigma(u).
\]
由于 Hilbert 变换在 $L^p$, $1<p<\infty$, 上有界, 我们证明了如下结论.

\begin{Corollary}
\label{cor:ch4-odd-kernel-lp}
若 $\Omega$ 是 $S^{n-1}$ 上的可积奇函数, 则 \eqref{eq:ch4-singular-integral-definition} 定义的算子 $T$ 在 $L^p(\mathbb R^n)$, $1<p<\infty$, 上有界.
\end{Corollary}

在推论~\ref{cor:ch4-odd-kernel-lp} 中, 我们把 $Tf$ 隐含地定义成 $L^p$ 范数下的极限. 与 Hilbert 变换一样, 还可以证明 \eqref{eq:ch4-singular-integral-definition} 对几乎处处的 $x$ 成立. 定义相应的极大截断算子
\begin{equation}
\label{eq:ch4-maximal-truncated-singular-integral}
 T^*f(x)=\sup_{\varepsilon>0}\left|
 \int_{|y|>\varepsilon}\frac{\Omega(y')}{|y|^n}f(x-y)\,dy\right|.
\end{equation}

\hypertarget{chapter-04-page-088}{}
用与上面相同的论证可得
\[
 T^*f(x)\le\frac\pi2\int_{S^{n-1}}|\Omega(u)|H_u^*f(x)\,d\sigma(u),
\]
其中 $H_u^*$ 是由极大 Hilbert 变换得到的方向算子. 因而由命题~\ref{prop:ch4-directional-operator-average} 和定理~\ref{thm:ch3-maximal-hilbert-bounds} 得到下述结论.

\begin{Corollary}
\label{cor:ch4-maximal-truncation-lp}
在与推论~\ref{cor:ch4-odd-kernel-lp} 相同的假设下, \eqref{eq:ch4-maximal-truncated-singular-integral} 定义的 $T^*$ 是强型 $(p,p)$ 的, $1<p<\infty$. 特别地, 对 $f\in L^p$, 极限 \eqref{eq:ch4-singular-integral-definition} 对几乎处处的 $x\in\mathbb R^n$ 成立.
\end{Corollary}

具有奇核的一族重要算子是 Riesz 变换:
\begin{equation}
\label{eq:ch4-riesz-transform-definition}
 R_jf(x)=c_n\operatorname{p.v.}\int_{\mathbb R^n}
 \frac{y_j}{|y|^{n+1}}f(x-y)\,dy,
 \qquad 1\le j\le n,
\end{equation}
其中
\[
 c_n=\Gamma\left(\frac{n+1}{2}\right)\pi^{-(n+1)/2}.
\]
这个常数选取为使得对 $f\in L^2$ 有
\begin{equation}
\label{eq:ch4-riesz-transform-multiplier}
 \widehat{R_jf}(\xi)=-i\frac{\xi_j}{|\xi|}\widehat f(\xi).
\end{equation}
进而
\begin{equation}
\label{eq:ch4-riesz-square-identity}
 \sum_{j=1}^nR_j^2=-I,
\end{equation}
其中 $I$ 是 $L^2$ 上的恒等算子. 由于 Schwartz 函数在 $L^p$, $p>1$, 中稠密, 这个恒等式实际上在每个 $L^p$ 空间中都成立.

为证明 \eqref{eq:ch4-riesz-transform-multiplier}, 可以直接使用定理~\ref{thm:ch4-fourier-principal-value-kernel}. 也可以注意到, 在分布意义下
\[
 \frac{\partial}{\partial x_j}|x|^{1-n}
 =(1-n)\operatorname{p.v.}\frac{x_j}{|x|^{n+1}},
\]
再应用 \eqref{eq:ch4-fourier-power-kernel} 与 Fourier 变换的微分性质.
